\documentclass[12pt,a4paper]{article}

\usepackage[T2A]{fontenc}
\usepackage[utf8]{inputenc}
\usepackage[russian]{babel}
\usepackage{amsmath,amssymb,amsthm}
\usepackage{enumitem}
\usepackage{listings}
\usepackage{xcolor}
\usepackage{geometry}
\usepackage{array}
\usepackage{longtable}
\usepackage{hyperref}

\geometry{left=2cm,right=2cm,top=2cm,bottom=2cm}

\lstdefinestyle{code}{
    basicstyle=\ttfamily\small,
    keywordstyle=\color{blue},
    commentstyle=\color{gray},
    stringstyle=\color{red!60!black},
    numbers=left,
    numberstyle=\tiny,
    stepnumber=1,
    numbersep=8pt,
    frame=single,
    breaklines=true,
    tabsize=4,
    showstringspaces=false,
    inputencoding=utf8,
    extendedchars=true,
    literate=
        {А}{{\CYRA}}1
        {Б}{{\CYRB}}1
        {В}{{\CYRV}}1
        {Г}{{\CYRG}}1
        {Д}{{\CYRD}}1
        {Е}{{\CYRE}}1
        {Ё}{{\CYRYO}}1
        {Ж}{{\CYRZH}}1
        {З}{{\CYRZ}}1
        {И}{{\CYRI}}1
        {Й}{{\CYRISHRT}}1
        {К}{{\CYRK}}1
        {Л}{{\CYRL}}1
        {М}{{\CYRM}}1
        {Н}{{\CYRN}}1
        {О}{{\CYRO}}1
        {П}{{\CYRP}}1
        {Р}{{\CYRR}}1
        {С}{{\CYRS}}1
        {Т}{{\CYRT}}1
        {У}{{\CYRU}}1
        {Ф}{{\CYRF}}1
        {Х}{{\CYRH}}1
        {Ц}{{\CYRC}}1
        {Ч}{{\CYRCH}}1
        {Ш}{{\CYRSH}}1
        {Щ}{{\CYRSHCH}}1
        {Ъ}{{\CYRHRDSN}}1
        {Ы}{{\CYRERY}}1
        {Ь}{{\CYRSFTSN}}1
        {Э}{{\CYREREV}}1
        {Ю}{{\CYRYU}}1
        {Я}{{\CYRYA}}1
        {а}{{\cyra}}1
        {б}{{\cyrb}}1
        {в}{{\cyrv}}1
        {г}{{\cyrg}}1
        {д}{{\cyrd}}1
        {е}{{\cyre}}1
        {ё}{{\cyryo}}1
        {ж}{{\cyrzh}}1
        {з}{{\cyrz}}1
        {и}{{\cyri}}1
        {й}{{\cyrishrt}}1
        {к}{{\cyrk}}1
        {л}{{\cyrl}}1
        {м}{{\cyrm}}1
        {н}{{\cyrn}}1
        {о}{{\cyro}}1
        {п}{{\cyrp}}1
        {р}{{\cyrr}}1
        {с}{{\cyrs}}1
        {т}{{\cyrt}}1
        {у}{{\cyru}}1
        {ф}{{\cyrf}}1
        {х}{{\cyrh}}1
        {ц}{{\cyrc}}1
        {ч}{{\cyrch}}1
        {ш}{{\cyrsh}}1
        {щ}{{\cyrshch}}1
        {ъ}{{\cyrhrdsn}}1
        {ы}{{\cyrery}}1
        {ь}{{\cyrsftsn}}1
        {э}{{\cyrerev}}1
        {ю}{{\cyryu}}1
        {я}{{\cyrya}}1
}

\lstset{style=code}

\title{Билет №44 \\ \small Объектно-ориентированное программирование. Шаблоны проектирования, их применение. Классификация шаблонов проектирования. Примеры шаблонов проектирования. (ИТМО) \\ Объектно-ориентированное программирование. Классы и объекты, наследование, интерфейсы. Понятие об объектном окружении. Рефлексия. Библиотеки классов. Средства обработки объектов (контейнеры, итераторы, сопроцедуры). (СпБГУ) \\ Основные принципы объектно-ориентированного программирования. (СпБПУ)}
\author{Сериков Владислав Александрович}
\date{13 мая 2026}

\begin{document}

\maketitle
\tableofcontents
\newpage

\section{Объектно-ориентированное программирование}

\subsection{Общая идея ООП}

\textbf{Объектно-ориентированное программирование} --- это парадигма программирования, в которой программа рассматривается как совокупность взаимодействующих объектов.

\textbf{Объект} объединяет:
\begin{itemize}
    \item \textbf{состояние} --- данные, хранящиеся в полях;
    \item \textbf{поведение} --- операции, реализованные методами;
    \item \textbf{идентичность} --- способность отличать один объект от другого даже при одинаковом состоянии.
\end{itemize}

Например, объект \texttt{bankAccount} может иметь состояние:
\[
\texttt{balance = 1000}
\]
и поведение:
\[
\texttt{deposit()}, \quad \texttt{withdraw()}, \quad \texttt{getBalance()}.
\]

ООП удобно применять там, где предметная область естественно описывается сущностями и их взаимодействием: банковские системы, графические интерфейсы, игры, информационные системы, компиляторы, системы моделирования.

\subsection{Класс и объект}

\textbf{Класс} --- это описание множества объектов с одинаковой структурой и поведением.

\textbf{Объект} --- это экземпляр класса, созданный во время выполнения программы.

\begin{lstlisting}[language=Java,caption={Пример класса и объекта}]
class Point {
    private int x;
    private int y;

    public Point(int x, int y) {
        this.x = x;
        this.y = y;
    }

    public void move(int dx, int dy) {
        x += dx;
        y += dy;
    }

    public String toString() {
        return "(" + x + ", " + y + ")";
    }
}

public class Main {
    public static void main(String[] args) {
        Point p = new Point(2, 3);
        p.move(1, -1);
        System.out.println(p); // (3, 2)
    }
}
\end{lstlisting}

Здесь \texttt{Point} --- класс, а \texttt{p} --- объект.

\subsection{Атрибуты, методы и конструкторы}

\textbf{Атрибуты} или \textbf{поля} класса задают состояние объекта.

\textbf{Методы} описывают действия, которые объект может выполнять.

\textbf{Конструктор} --- специальный метод, вызываемый при создании объекта. Его задача --- привести объект в корректное начальное состояние.

\begin{lstlisting}[language=Java,caption={Конструктор и методы}]
class User {
    private String name;
    private int age;

    public User(String name, int age) {
        if (age < 0) {
            throw new IllegalArgumentException("age must be non-negative");
        }
        this.name = name;
        this.age = age;
    }

    public String getName() {
        return name;
    }
}
\end{lstlisting}

\subsection{Основные принципы ООП}

Классически выделяют четыре основных принципа ООП:
\begin{enumerate}
    \item абстракция;
    \item инкапсуляция;
    \item наследование;
    \item полиморфизм.
\end{enumerate}

Иногда дополнительно выделяют:
\begin{itemize}
    \item композицию;
    \item динамическое связывание;
    \item посылку сообщений;
    \item модульность;
    \item повторное использование кода.
\end{itemize}

\section{Абстракция}

\subsection{Понятие абстракции}

\textbf{Абстракция} --- это выделение существенных свойств объекта и игнорирование несущественных деталей реализации.

В программировании абстракция позволяет работать с объектом через его интерфейс, не зная внутреннего устройства.

Например, при использовании списка важно знать, что у него есть операции:
\[
\texttt{add}, \quad \texttt{remove}, \quad \texttt{get}, \quad \texttt{size},
\]
но не обязательно знать, хранится ли список в массиве, связном списке или дереве.

\subsection{Пример абстракции}

\begin{lstlisting}[language=Java,caption={Абстрактный тип данных}]
interface Stack<T> {
    void push(T value);
    T pop();
    T peek();
    boolean isEmpty();
}
\end{lstlisting}

Интерфейс \texttt{Stack} описывает поведение стека, но не задаёт конкретный способ хранения элементов.

\subsection{Абстрактные классы}

\textbf{Абстрактный класс} --- класс, который может содержать как реализованные методы, так и абстрактные методы без реализации.

\begin{lstlisting}[language=Java,caption={Абстрактный класс}]
abstract class Shape {
    public abstract double area();

    public void printArea() {
        System.out.println("Area = " + area());
    }
}

class Circle extends Shape {
    private double r;

    public Circle(double r) {
        this.r = r;
    }

    public double area() {
        return Math.PI * r * r;
    }
}
\end{lstlisting}

Абстрактный класс полезен, когда разные классы имеют общую часть реализации.

\section{Инкапсуляция}

\subsection{Понятие инкапсуляции}

\textbf{Инкапсуляция} --- это объединение данных и методов, работающих с этими данными, внутри одной программной единицы, а также сокрытие внутреннего состояния объекта от внешнего кода.

Идея инкапсуляции выражается следующим правилом:

\[
\text{внешний код должен зависеть от интерфейса объекта, а не от его реализации.}
\]

\subsection{Модификаторы доступа}

В объектно-ориентированных языках обычно существуют средства управления доступом:
\begin{itemize}
    \item \texttt{public} --- доступно всем;
    \item \texttt{private} --- доступно только внутри класса;
    \item \texttt{protected} --- доступно внутри класса, его наследников и иногда внутри пакета;
    \item package-private или internal-доступ --- доступ в пределах модуля или пакета.
\end{itemize}

\subsection{Пример инкапсуляции}

\begin{lstlisting}[language=Java,caption={Инкапсуляция состояния}]
class BankAccount {
    private int balance;

    public BankAccount(int initialBalance) {
        if (initialBalance < 0) {
            throw new IllegalArgumentException("Negative balance");
        }
        this.balance = initialBalance;
    }

    public void deposit(int amount) {
        if (amount <= 0) {
            throw new IllegalArgumentException("Amount must be positive");
        }
        balance += amount;
    }

    public void withdraw(int amount) {
        if (amount <= 0 || amount > balance) {
            throw new IllegalArgumentException("Invalid amount");
        }
        balance -= amount;
    }

    public int getBalance() {
        return balance;
    }
}
\end{lstlisting}

Поле \texttt{balance} закрыто от прямого изменения. Это позволяет гарантировать инвариант:
\[
\texttt{balance} \geq 0.
\]

\subsection{Инвариант класса}

\textbf{Инвариант класса} --- логическое условие, которое должно быть истинно для каждого корректного объекта класса после завершения конструктора и после выполнения каждого публичного метода.

Для класса \texttt{BankAccount} инвариант:
\[
balance \geq 0.
\]

Инкапсуляция помогает поддерживать инварианты, потому что запрещает внешнему коду напрямую менять внутренние поля объекта.

\section{Наследование}

\subsection{Понятие наследования}

\textbf{Наследование} --- механизм, позволяющий определить новый класс на основе уже существующего.

Существующий класс называется:
\begin{itemize}
    \item базовым классом;
    \item родительским классом;
    \item суперклассом.
\end{itemize}

Новый класс называется:
\begin{itemize}
    \item производным классом;
    \item дочерним классом;
    \item подклассом.
\end{itemize}

Наследование выражает отношение:
\[
\text{``является''} \quad \text{или} \quad is\text{-}a.
\]

Например:
\[
\texttt{Circle is a Shape}.
\]

\subsection{Пример наследования}

\begin{lstlisting}[language=Java,caption={Наследование}]
class Animal {
    public void eat() {
        System.out.println("Animal eats");
    }
}

class Dog extends Animal {
    public void bark() {
        System.out.println("Dog barks");
    }
}

public class Main {
    public static void main(String[] args) {
        Dog dog = new Dog();
        dog.eat();
        dog.bark();
    }
}
\end{lstlisting}

Класс \texttt{Dog} наследует метод \texttt{eat} от класса \texttt{Animal}.

\subsection{Переопределение методов}

\textbf{Переопределение метода} --- определение в подклассе метода с той же сигнатурой, что и у метода базового класса.

\begin{lstlisting}[language=Java,caption={Переопределение метода}]
class Animal {
    public void speak() {
        System.out.println("Some sound");
    }
}

class Cat extends Animal {
    @Override
    public void speak() {
        System.out.println("Meow");
    }
}
\end{lstlisting}

\subsection{Наследование реализации и наследование интерфейса}

Следует различать:
\begin{itemize}
    \item \textbf{наследование реализации} --- подкласс получает готовый код базового класса;
    \item \textbf{наследование интерфейса} --- класс обязуется поддерживать некоторый набор операций.
\end{itemize}

Наследование реализации удобно для повторного использования кода, но может создавать сильную зависимость между классами.

Наследование интерфейса обычно является более гибким, так как оно связывает классы через контракт, а не через реализацию.

\subsection{Проблемы наследования}

Наследование может приводить к следующим проблемам:
\begin{itemize}
    \item сильная связанность базового и производного классов;
    \item хрупкость иерархий;
    \item нарушение инкапсуляции;
    \item сложность множественного наследования;
    \item трудности при изменении базового класса.
\end{itemize}

Поэтому часто используют принцип:

\[
\text{предпочитать композицию наследованию.}
\]

\section{Композиция и агрегация}

\subsection{Композиция}

\textbf{Композиция} --- отношение между объектами вида ``часть--целое'', при котором один объект содержит другой и управляет его жизненным циклом.

Например, объект \texttt{Car} может содержать объект \texttt{Engine}.

\begin{lstlisting}[language=Java,caption={Композиция}]
class Engine {
    public void start() {
        System.out.println("Engine started");
    }
}

class Car {
    private Engine engine = new Engine();

    public void start() {
        engine.start();
    }
}
\end{lstlisting}

Композиция выражает отношение:
\[
\text{has-a}.
\]

То есть:
\[
\texttt{Car has an Engine}.
\]

\subsection{Агрегация}

\textbf{Агрегация} --- более слабое отношение ``часть--целое'', при котором объект использует другой объект, но не обязательно управляет его жизненным циклом.

Пример:
\[
\texttt{Team has Players}.
\]

Игрок может существовать независимо от команды.

\section{Полиморфизм}

\subsection{Понятие полиморфизма}

\textbf{Полиморфизм} --- возможность использовать объекты разных классов единообразно через общий интерфейс.

В ООП чаще всего рассматривают:
\begin{itemize}
    \item полиморфизм подтипов;
    \item параметрический полиморфизм;
    \item ad-hoc полиморфизм.
\end{itemize}

\subsection{Полиморфизм подтипов}

\textbf{Полиморфизм подтипов} означает, что объект производного класса может использоваться там, где ожидается объект базового класса или интерфейса.

\begin{lstlisting}[language=Java,caption={Полиморфизм подтипов}]
interface Shape {
    double area();
}

class Rectangle implements Shape {
    private double width;
    private double height;

    public Rectangle(double width, double height) {
        this.width = width;
        this.height = height;
    }

    public double area() {
        return width * height;
    }
}

class Circle implements Shape {
    private double r;

    public Circle(double r) {
        this.r = r;
    }

    public double area() {
        return Math.PI * r * r;
    }
}

public class Main {
    public static void printArea(Shape shape) {
        System.out.println(shape.area());
    }

    public static void main(String[] args) {
        printArea(new Rectangle(3, 4));
        printArea(new Circle(2));
    }
}
\end{lstlisting}

Функция \texttt{printArea} работает с любым объектом, реализующим интерфейс \texttt{Shape}.

\subsection{Динамическое связывание}

\textbf{Динамическое связывание} --- выбор конкретной реализации метода во время выполнения программы, исходя из фактического типа объекта.

Если переменная имеет тип \texttt{Shape}, но фактически ссылается на объект \texttt{Circle}, то будет вызван метод \texttt{area} класса \texttt{Circle}.

\[
\texttt{Shape s = new Circle(2);}
\]
\[
\texttt{s.area();}
\]

Компилятор знает только, что \texttt{s} имеет тип \texttt{Shape}, но во время выполнения выбирается реализация \texttt{Circle.area()}.

\subsection{Параметрический полиморфизм}

\textbf{Параметрический полиморфизм} реализуется с помощью обобщений или шаблонов.

\begin{lstlisting}[language=Java,caption={Параметрический полиморфизм}]
class Box<T> {
    private T value;

    public Box(T value) {
        this.value = value;
    }

    public T get() {
        return value;
    }
}
\end{lstlisting}

Класс \texttt{Box<T>} может работать с разными типами:
\[
\texttt{Box<Integer>}, \quad \texttt{Box<String>}.
\]

\subsection{Ad-hoc полиморфизм}

\textbf{Ad-hoc полиморфизм} обычно реализуется перегрузкой функций и операторов.

\begin{lstlisting}[language=Java,caption={Перегрузка методов}]
class Printer {
    public void print(int x) {
        System.out.println(x);
    }

    public void print(String s) {
        System.out.println(s);
    }
}
\end{lstlisting}

Методы имеют одно имя, но разные параметры.

\section{Интерфейсы}

\subsection{Понятие интерфейса}

\textbf{Интерфейс} --- это описание контракта, который должен выполнять класс.

Интерфейс задаёт, какие операции доступны, но обычно не навязывает конкретную реализацию.

\begin{lstlisting}[language=Java,caption={Интерфейс}]
interface Repository<T> {
    void save(T object);
    T findById(int id);
}
\end{lstlisting}

Класс, реализующий интерфейс, обязан предоставить реализации его методов.

\begin{lstlisting}[language=Java,caption={Реализация интерфейса}]
class UserRepository implements Repository<User> {
    public void save(User user) {
        // сохранение пользователя
    }

    public User findById(int id) {
        // поиск пользователя
        return null;
    }
}
\end{lstlisting}

\subsection{Роль интерфейсов}

Интерфейсы используются для:
\begin{itemize}
    \item задания контрактов;
    \item ослабления связанности между компонентами;
    \item реализации полиморфизма;
    \item поддержки тестирования;
    \item замены конкретных реализаций без изменения клиентского кода.
\end{itemize}

\subsection{Интерфейс и абстрактный класс}

\begin{longtable}{|p{5cm}|p{5cm}|p{5cm}|}
\hline
\textbf{Критерий} & \textbf{Интерфейс} & \textbf{Абстрактный класс} \\
\hline
Назначение & Контракт поведения & Частичная реализация общей сущности \\
\hline
Состояние & Обычно не содержит изменяемого состояния объекта & Может содержать поля \\
\hline
Множественность & Класс обычно может реализовать несколько интерфейсов & Во многих языках класс наследует только один класс \\
\hline
Связь с реализацией & Слабая & Более сильная \\
\hline
Применение & Когда важен набор операций & Когда есть общая логика и состояние \\
\hline
\end{longtable}

\section{Объектное окружение}

\subsection{Понятие объектного окружения}

\textbf{Объектное окружение} --- совокупность средств языка, среды выполнения и библиотек, которые обеспечивают создание, хранение, взаимодействие, уничтожение и обработку объектов.

В объектное окружение входят:
\begin{itemize}
    \item система типов;
    \item механизм создания объектов;
    \item механизм вызова методов;
    \item система управления памятью;
    \item стандартная библиотека классов;
    \item контейнеры;
    \item итераторы;
    \item исключения;
    \item механизмы рефлексии;
    \item средства сериализации;
    \item механизмы параллельного и асинхронного выполнения;
    \item рантайм-среда.
\end{itemize}

Например, в Java объектное окружение включает:
\begin{itemize}
    \item виртуальную машину JVM;
    \item загрузчики классов;
    \item сборщик мусора;
    \item стандартную библиотеку;
    \item механизм рефлексии;
    \item систему исключений;
    \item коллекции;
    \item потоки и средства конкурентного программирования.
\end{itemize}

\subsection{Жизненный цикл объекта}

Типичный жизненный цикл объекта:
\begin{enumerate}
    \item выделение памяти;
    \item вызов конструктора;
    \item использование объекта;
    \item потеря достижимости объекта;
    \item освобождение ресурсов вручную или сборщиком мусора.
\end{enumerate}

\subsection{Схема взаимодействия объектов}

Пусть есть объект \texttt{Client}, объект \texttt{Service} и объект \texttt{Repository}.

\[
\texttt{Client} \rightarrow \texttt{Service} \rightarrow \texttt{Repository}.
\]

Клиент вызывает сервис, сервис выполняет бизнес-логику и обращается к репозиторию.

\begin{lstlisting}[language=Java,caption={Простое объектное окружение}]
interface Repository {
    String findNameById(int id);
}

class MemoryRepository implements Repository {
    public String findNameById(int id) {
        return "Alice";
    }
}

class UserService {
    private Repository repository;

    public UserService(Repository repository) {
        this.repository = repository;
    }

    public String getUserName(int id) {
        return repository.findNameById(id);
    }
}

public class Main {
    public static void main(String[] args) {
        Repository repository = new MemoryRepository();
        UserService service = new UserService(repository);
        System.out.println(service.getUserName(1));
    }
}
\end{lstlisting}

Здесь объект \texttt{UserService} не зависит от конкретного класса \texttt{MemoryRepository}, а зависит только от интерфейса \texttt{Repository}.

\section{Рефлексия}

\subsection{Понятие рефлексии}

\textbf{Рефлексия} --- механизм, позволяющий программе исследовать и иногда изменять свою структуру и поведение во время выполнения.

С помощью рефлексии можно:
\begin{itemize}
    \item получить информацию о классе объекта;
    \item узнать список методов и полей;
    \item вызвать метод по имени;
    \item создать объект по имени класса;
    \item получить аннотации или метаданные;
    \item изменить доступность полей и методов;
    \item построить универсальные фреймворки.
\end{itemize}

\subsection{Пример рефлексии}

\begin{lstlisting}[language=Java,caption={Получение информации о классе}]
import java.lang.reflect.Method;

class Person {
    public void sayHello() {
        System.out.println("Hello");
    }
}

public class Main {
    public static void main(String[] args) {
        Class<?> clazz = Person.class;

        for (Method method : clazz.getDeclaredMethods()) {
            System.out.println(method.getName());
        }
    }
}
\end{lstlisting}

Результат:
\[
\texttt{sayHello}
\]

\subsection{Создание объекта через рефлексию}

\begin{lstlisting}[language=Java,caption={Создание объекта через рефлексию}]
class Person {
    public Person() {}

    public void sayHello() {
        System.out.println("Hello");
    }
}

public class Main {
    public static void main(String[] args) throws Exception {
        Class<?> clazz = Class.forName("Person");
        Object obj = clazz.getDeclaredConstructor().newInstance();

        Person person = (Person) obj;
        person.sayHello();
    }
}
\end{lstlisting}

\subsection{Применение рефлексии}

Рефлексия применяется:
\begin{itemize}
    \item в системах сериализации;
    \item в ORM-фреймворках;
    \item в контейнерах зависимостей;
    \item в тестовых фреймворках;
    \item в плагинных архитектурах;
    \item в средах разработки;
    \item в отладчиках.
\end{itemize}

\subsection{Недостатки рефлексии}

Недостатки:
\begin{itemize}
    \item снижение производительности;
    \item потеря статической проверки типов;
    \item возможность нарушения инкапсуляции;
    \item усложнение сопровождения;
    \item потенциальные проблемы безопасности.
\end{itemize}

\section{Библиотеки классов}

\subsection{Понятие библиотеки классов}

\textbf{Библиотека классов} --- набор готовых классов, интерфейсов и функций, предназначенных для повторного использования.

Библиотеки классов предоставляют:
\begin{itemize}
    \item структуры данных;
    \item алгоритмы;
    \item средства ввода-вывода;
    \item работу с файлами;
    \item сетевое взаимодействие;
    \item графический интерфейс;
    \item средства многопоточности;
    \item работу с датой и временем;
    \item математические функции;
    \item механизмы сериализации;
    \item средства тестирования.
\end{itemize}

\subsection{Примеры библиотек классов}

Примеры:
\begin{itemize}
    \item Java Standard Library;
    \item .NET Base Class Library;
    \item C++ Standard Library;
    \item Qt;
    \item Boost;
    \item Python Standard Library.
\end{itemize}

\subsection{Польза библиотек классов}

Библиотеки классов:
\begin{itemize}
    \item уменьшают объём кода;
    \item повышают надёжность;
    \item дают проверенные реализации структур данных и алгоритмов;
    \item стандартизируют подходы;
    \item ускоряют разработку;
    \item упрощают сопровождение.
\end{itemize}

\section{Средства обработки объектов}

\subsection{Контейнеры}

\textbf{Контейнер} --- объект, предназначенный для хранения других объектов.

Основные виды контейнеров:
\begin{itemize}
    \item массив;
    \item список;
    \item стек;
    \item очередь;
    \item дек;
    \item множество;
    \item ассоциативный массив;
    \item дерево;
    \item граф.
\end{itemize}

\subsection{Список}

\textbf{Список} --- последовательный контейнер, поддерживающий хранение элементов в определённом порядке.

\begin{lstlisting}[language=Java,caption={Список}]
import java.util.*;

public class Main {
    public static void main(String[] args) {
        List<String> names = new ArrayList<>();

        names.add("Alice");
        names.add("Bob");

        System.out.println(names.get(0)); // Alice
    }
}
\end{lstlisting}

\subsection{Множество}

\textbf{Множество} --- контейнер, не допускающий повторяющихся элементов.

\begin{lstlisting}[language=Java,caption={Множество}]
import java.util.*;

public class Main {
    public static void main(String[] args) {
        Set<Integer> numbers = new HashSet<>();

        numbers.add(1);
        numbers.add(1);
        numbers.add(2);

        System.out.println(numbers.size()); // 2
    }
}
\end{lstlisting}

\subsection{Ассоциативный массив}

\textbf{Ассоциативный массив} или \textbf{словарь} хранит пары:
\[
\text{ключ} \rightarrow \text{значение}.
\]

\begin{lstlisting}[language=Java,caption={Ассоциативный массив}]
import java.util.*;

public class Main {
    public static void main(String[] args) {
        Map<String, Integer> ages = new HashMap<>();

        ages.put("Alice", 20);
        ages.put("Bob", 25);

        System.out.println(ages.get("Alice")); // 20
    }
}
\end{lstlisting}

\subsection{Итераторы}

\textbf{Итератор} --- объект, предоставляющий последовательный доступ к элементам контейнера без раскрытия его внутреннего устройства.

Обычно итератор поддерживает операции:
\begin{itemize}
    \item проверить наличие следующего элемента;
    \item получить следующий элемент.
\end{itemize}

\subsection{Алгоритм обхода контейнера с помощью итератора}

Пусть есть контейнер $C$.

\textbf{Шаги алгоритма:}
\begin{enumerate}
    \item Получить итератор $it$ для контейнера $C$.
    \item Пока в контейнере есть следующий элемент:
    \begin{enumerate}
        \item получить текущий элемент;
        \item обработать текущий элемент;
        \item перейти к следующему элементу.
    \end{enumerate}
    \item Завершить обход.
\end{enumerate}

\textbf{Иллюстрация:}

\[
[10, 20, 30]
\]

\[
it \rightarrow 10 \quad \Rightarrow \quad it \rightarrow 20 \quad \Rightarrow \quad it \rightarrow 30.
\]

\begin{lstlisting}[language=Java,caption={Обход с помощью итератора}]
import java.util.*;

public class Main {
    public static void main(String[] args) {
        List<Integer> values = Arrays.asList(10, 20, 30);

        Iterator<Integer> it = values.iterator();

        while (it.hasNext()) {
            int x = it.next();
            System.out.println(x);
        }
    }
}
\end{lstlisting}

\subsection{Сопроцедуры и корутины}

\textbf{Сопроцедура} или \textbf{корутина} --- программный компонент, который может приостанавливать своё выполнение и затем продолжать его с сохранённого состояния.

Обычная процедура имеет схему:

\[
\text{вызов} \rightarrow \text{выполнение} \rightarrow \text{возврат}.
\]

Корутина имеет схему:

\[
\text{вызов} \rightarrow \text{выполнение} \rightarrow \text{приостановка} \rightarrow \text{возобновление} \rightarrow \dots
\]

\subsection{Генераторы как частный случай корутин}

Генератор выдаёт значения по одному, сохраняя своё состояние между вызовами.

\begin{lstlisting}[language=Python,caption={Генератор как пример корутины}]
def numbers():
    yield 1
    yield 2
    yield 3

for x in numbers():
    print(x)
\end{lstlisting}

\textbf{Шаги выполнения:}
\begin{enumerate}
    \item Вызов \texttt{numbers()} создаёт объект-генератор.
    \item Первый вызов итерации выполняет код до \texttt{yield 1}.
    \item Значение $1$ возвращается вызывающему коду.
    \item Состояние функции сохраняется.
    \item Следующая итерация продолжает выполнение после \texttt{yield 1}.
    \item Аналогично возвращаются $2$ и $3$.
    \item После завершения функции итерация заканчивается.
\end{enumerate}

\textbf{Иллюстрация:}

\[
start \rightarrow yield\ 1 \rightarrow pause
\]

\[
resume \rightarrow yield\ 2 \rightarrow pause
\]

\[
resume \rightarrow yield\ 3 \rightarrow finish.
\]

\section{Шаблоны проектирования}

\subsection{Понятие шаблона проектирования}

\textbf{Шаблон проектирования} --- это обобщённое повторно используемое решение типичной задачи проектирования программного обеспечения.

Шаблон не является готовым фрагментом кода. Он описывает:
\begin{itemize}
    \item проблему;
    \item контекст её возникновения;
    \item структуру решения;
    \item участников;
    \item последствия применения.
\end{itemize}

Шаблоны проектирования позволяют использовать накопленный опыт разработки и дают общий язык общения между разработчиками.

\subsection{Зачем нужны шаблоны проектирования}

Шаблоны применяются для:
\begin{itemize}
    \item повышения повторного использования решений;
    \item уменьшения связанности между классами;
    \item повышения гибкости архитектуры;
    \item упрощения сопровождения;
    \item стандартизации проектных решений;
    \item документирования архитектурных идей.
\end{itemize}

\subsection{Опасность чрезмерного применения}

Шаблоны не следует применять механически.

Недостатки чрезмерного использования:
\begin{itemize}
    \item усложнение архитектуры;
    \item появление лишних классов;
    \item ухудшение читаемости;
    \item снижение производительности;
    \item переинжиниринг простых задач.
\end{itemize}

Шаблон оправдан, если он решает реальную проблему изменяемости, расширяемости или повторного использования.

\section{Классификация шаблонов проектирования}

Классическая классификация шаблонов проектирования делит их на три группы:
\begin{enumerate}
    \item порождающие;
    \item структурные;
    \item поведенческие.
\end{enumerate}

\subsection{Порождающие шаблоны}

\textbf{Порождающие шаблоны} управляют процессом создания объектов.

Они помогают:
\begin{itemize}
    \item скрыть конкретные классы создаваемых объектов;
    \item отделить создание объекта от его использования;
    \item обеспечить гибкую настройку процесса создания.
\end{itemize}

К порождающим шаблонам относятся:
\begin{itemize}
    \item Singleton;
    \item Factory Method;
    \item Abstract Factory;
    \item Builder;
    \item Prototype.
\end{itemize}

\subsection{Структурные шаблоны}

\textbf{Структурные шаблоны} описывают способы объединения классов и объектов в более крупные структуры.

Они помогают:
\begin{itemize}
    \item адаптировать интерфейсы;
    \item строить иерархии объектов;
    \item добавлять функциональность без изменения исходного класса;
    \item уменьшать связанность.
\end{itemize}

К структурным шаблонам относятся:
\begin{itemize}
    \item Adapter;
    \item Bridge;
    \item Composite;
    \item Decorator;
    \item Facade;
    \item Flyweight;
    \item Proxy.
\end{itemize}

\subsection{Поведенческие шаблоны}

\textbf{Поведенческие шаблоны} описывают взаимодействие объектов и распределение ответственности между ними.

Они помогают:
\begin{itemize}
    \item организовывать обмен сообщениями;
    \item изменять поведение объекта во время выполнения;
    \item отделять алгоритмы от объектов;
    \item уменьшать связанность участников взаимодействия.
\end{itemize}

К поведенческим шаблонам относятся:
\begin{itemize}
    \item Chain of Responsibility;
    \item Command;
    \item Interpreter;
    \item Iterator;
    \item Mediator;
    \item Memento;
    \item Observer;
    \item State;
    \item Strategy;
    \item Template Method;
    \item Visitor.
\end{itemize}

\section{Примеры шаблонов проектирования}

\subsection{Singleton}

\subsubsection{Назначение}

\textbf{Singleton} гарантирует, что у класса существует только один экземпляр, и предоставляет глобальную точку доступа к нему.

\subsubsection{Когда применять}

Шаблон применяют, когда:
\begin{itemize}
    \item в системе должен быть ровно один объект данного класса;
    \item нужен централизованный доступ к общему ресурсу;
    \item создание нескольких экземпляров нарушает логику программы.
\end{itemize}

Примеры:
\begin{itemize}
    \item конфигурация приложения;
    \item журналирование;
    \item пул соединений;
    \item менеджер ресурсов.
\end{itemize}

\subsubsection{Минимальный пример}

\begin{lstlisting}[language=Java,caption={Singleton}]
class Logger {
    private static final Logger INSTANCE = new Logger();

    private Logger() {}

    public static Logger getInstance() {
        return INSTANCE;
    }

    public void log(String message) {
        System.out.println(message);
    }
}

public class Main {
    public static void main(String[] args) {
        Logger logger = Logger.getInstance();
        logger.log("Application started");
    }
}
\end{lstlisting}

\subsubsection{Достоинства и недостатки}

Достоинства:
\begin{itemize}
    \item контроль числа экземпляров;
    \item единая точка доступа.
\end{itemize}

Недостатки:
\begin{itemize}
    \item глобальное состояние;
    \item усложнение тестирования;
    \item скрытые зависимости;
    \item возможное нарушение принципа единственной ответственности.
\end{itemize}

\subsection{Factory Method}

\subsubsection{Назначение}

\textbf{Factory Method} определяет общий интерфейс создания объекта, но позволяет подклассам выбирать конкретный создаваемый класс.

\subsubsection{Идея}

Клиентский код работает с абстрактным продуктом, а создание конкретного продукта делегируется фабричному методу.

\[
\texttt{Creator} \rightarrow \texttt{factoryMethod()} \rightarrow \texttt{Product}.
\]

\subsubsection{Минимальный пример}

\begin{lstlisting}[language=Java,caption={Factory Method}]
interface Button {
    void render();
}

class WindowsButton implements Button {
    public void render() {
        System.out.println("Windows button");
    }
}

class HtmlButton implements Button {
    public void render() {
        System.out.println("HTML button");
    }
}

abstract class Dialog {
    public void renderDialog() {
        Button button = createButton();
        button.render();
    }

    protected abstract Button createButton();
}

class WindowsDialog extends Dialog {
    protected Button createButton() {
        return new WindowsButton();
    }
}

class WebDialog extends Dialog {
    protected Button createButton() {
        return new HtmlButton();
    }
}
\end{lstlisting}

\subsubsection{Алгоритм применения}

\begin{enumerate}
    \item Выделить общий интерфейс продукта.
    \item Создать конкретные реализации продукта.
    \item В базовом классе определить фабричный метод.
    \item В подклассах переопределить фабричный метод.
    \item Клиентский код должен работать с абстрактным продуктом.
\end{enumerate}

\subsection{Abstract Factory}

\subsubsection{Назначение}

\textbf{Abstract Factory} предоставляет интерфейс для создания семейств взаимосвязанных объектов без указания их конкретных классов.

\subsubsection{Пример задачи}

Нужно создавать элементы интерфейса для разных платформ:
\[
\texttt{WindowsButton}, \quad \texttt{WindowsCheckbox}
\]
и
\[
\texttt{MacButton}, \quad \texttt{MacCheckbox}.
\]

Важно, чтобы элементы одного семейства были совместимы.

\subsubsection{Минимальный пример}

\begin{lstlisting}[language=Java,caption={Abstract Factory}]
interface Button {
    void paint();
}

interface Checkbox {
    void paint();
}

class WindowsButton implements Button {
    public void paint() {
        System.out.println("Windows button");
    }
}

class WindowsCheckbox implements Checkbox {
    public void paint() {
        System.out.println("Windows checkbox");
    }
}

class MacButton implements Button {
    public void paint() {
        System.out.println("Mac button");
    }
}

class MacCheckbox implements Checkbox {
    public void paint() {
        System.out.println("Mac checkbox");
    }
}

interface GUIFactory {
    Button createButton();
    Checkbox createCheckbox();
}

class WindowsFactory implements GUIFactory {
    public Button createButton() {
        return new WindowsButton();
    }

    public Checkbox createCheckbox() {
        return new WindowsCheckbox();
    }
}

class MacFactory implements GUIFactory {
    public Button createButton() {
        return new MacButton();
    }

    public Checkbox createCheckbox() {
        return new MacCheckbox();
    }
}
\end{lstlisting}

\subsection{Builder}

\subsubsection{Назначение}

\textbf{Builder} отделяет конструирование сложного объекта от его представления.

Шаблон полезен, когда объект имеет много параметров или сложный процесс создания.

\subsubsection{Минимальный пример}

\begin{lstlisting}[language=Java,caption={Builder}]
class User {
    private String name;
    private int age;
    private String email;

    private User(Builder builder) {
        this.name = builder.name;
        this.age = builder.age;
        this.email = builder.email;
    }

    static class Builder {
        private String name;
        private int age;
        private String email;

        public Builder name(String name) {
            this.name = name;
            return this;
        }

        public Builder age(int age) {
            this.age = age;
            return this;
        }

        public Builder email(String email) {
            this.email = email;
            return this;
        }

        public User build() {
            return new User(this);
        }
    }
}

public class Main {
    public static void main(String[] args) {
        User user = new User.Builder()
            .name("Alice")
            .age(20)
            .email("alice@example.com")
            .build();
    }
}
\end{lstlisting}

\subsubsection{Шаги построения объекта}

\begin{enumerate}
    \item Создаётся объект-строитель.
    \item Последовательно задаются нужные параметры.
    \item Вызывается метод \texttt{build()}.
    \item Возвращается готовый объект.
\end{enumerate}

\subsection{Prototype}

\subsubsection{Назначение}

\textbf{Prototype} создаёт новые объекты путём копирования существующего объекта-прототипа.

\subsubsection{Когда применять}

Шаблон полезен, когда:
\begin{itemize}
    \item создание объекта дорогостоящее;
    \item объект имеет сложную начальную конфигурацию;
    \item нужно создавать объекты, похожие на уже существующие.
\end{itemize}

\subsubsection{Минимальный пример}

\begin{lstlisting}[language=Java,caption={Prototype}]
class Document implements Cloneable {
    private String text;

    public Document(String text) {
        this.text = text;
    }

    public Document clone() {
        return new Document(this.text);
    }
}
\end{lstlisting}

\subsection{Adapter}

\subsubsection{Назначение}

\textbf{Adapter} преобразует интерфейс одного класса в интерфейс, ожидаемый клиентом.

\subsubsection{Пример}

Есть старый класс:

\begin{lstlisting}[language=Java]
class OldPrinter {
    public void printOld(String text) {
        System.out.println(text);
    }
}
\end{lstlisting}

Клиент ожидает интерфейс:

\begin{lstlisting}[language=Java]
interface Printer {
    void print(String text);
}
\end{lstlisting}

Адаптер:

\begin{lstlisting}[language=Java,caption={Adapter}]
class PrinterAdapter implements Printer {
    private OldPrinter oldPrinter;

    public PrinterAdapter(OldPrinter oldPrinter) {
        this.oldPrinter = oldPrinter;
    }

    public void print(String text) {
        oldPrinter.printOld(text);
    }
}
\end{lstlisting}

\subsubsection{Иллюстрация}

\[
\texttt{Client} \rightarrow \texttt{Printer}
\]

\[
\texttt{PrinterAdapter} \rightarrow \texttt{OldPrinter}.
\]

Клиент не знает о существовании \texttt{OldPrinter}.

\subsection{Decorator}

\subsubsection{Назначение}

\textbf{Decorator} динамически добавляет объекту новую функциональность, не изменяя его класс.

\subsubsection{Минимальный пример}

\begin{lstlisting}[language=Java,caption={Decorator}]
interface Notifier {
    void send(String message);
}

class EmailNotifier implements Notifier {
    public void send(String message) {
        System.out.println("Email: " + message);
    }
}

class SmsDecorator implements Notifier {
    private Notifier notifier;

    public SmsDecorator(Notifier notifier) {
        this.notifier = notifier;
    }

    public void send(String message) {
        notifier.send(message);
        System.out.println("SMS: " + message);
    }
}

public class Main {
    public static void main(String[] args) {
        Notifier notifier = new SmsDecorator(new EmailNotifier());
        notifier.send("Hello");
    }
}
\end{lstlisting}

\subsubsection{Идея}

\[
\texttt{Decorator} \supset \texttt{Component}
\]

Декоратор содержит ссылку на объект того же интерфейса и дополняет его поведение.

\subsection{Facade}

\subsubsection{Назначение}

\textbf{Facade} предоставляет простой интерфейс к сложной подсистеме.

\subsubsection{Минимальный пример}

\begin{lstlisting}[language=Java,caption={Facade}]
class CPU {
    void freeze() {}
    void execute() {}
}

class Memory {
    void load() {}
}

class HardDrive {
    byte[] read() {
        return new byte[0];
    }
}

class ComputerFacade {
    private CPU cpu = new CPU();
    private Memory memory = new Memory();
    private HardDrive hardDrive = new HardDrive();

    public void start() {
        cpu.freeze();
        memory.load();
        hardDrive.read();
        cpu.execute();
    }
}
\end{lstlisting}

Клиент вызывает только:
\[
\texttt{computer.start();}
\]

и не работает напрямую со всеми частями подсистемы.

\subsection{Proxy}

\subsubsection{Назначение}

\textbf{Proxy} предоставляет объект-заместитель, контролирующий доступ к другому объекту.

\subsubsection{Применения}

\begin{itemize}
    \item ленивое создание объекта;
    \item контроль доступа;
    \item удалённый доступ;
    \item кеширование;
    \item логирование вызовов.
\end{itemize}

\subsubsection{Минимальный пример}

\begin{lstlisting}[language=Java,caption={Proxy}]
interface Image {
    void display();
}

class RealImage implements Image {
    private String filename;

    public RealImage(String filename) {
        this.filename = filename;
        loadFromDisk();
    }

    private void loadFromDisk() {
        System.out.println("Loading " + filename);
    }

    public void display() {
        System.out.println("Displaying " + filename);
    }
}

class ProxyImage implements Image {
    private String filename;
    private RealImage realImage;

    public ProxyImage(String filename) {
        this.filename = filename;
    }

    public void display() {
        if (realImage == null) {
            realImage = new RealImage(filename);
        }
        realImage.display();
    }
}
\end{lstlisting}

\subsection{Composite}

\subsubsection{Назначение}

\textbf{Composite} позволяет работать с отдельными объектами и группами объектов единообразно.

\subsubsection{Пример}

Файловая система:
\[
\texttt{File} \quad \text{и} \quad \texttt{Directory}
\]
могут иметь общий интерфейс \texttt{FileSystemNode}.

\begin{lstlisting}[language=Java,caption={Composite}]
import java.util.*;

interface Graphic {
    void draw();
}

class Circle implements Graphic {
    public void draw() {
        System.out.println("Circle");
    }
}

class Group implements Graphic {
    private List<Graphic> children = new ArrayList<>();

    public void add(Graphic graphic) {
        children.add(graphic);
    }

    public void draw() {
        for (Graphic graphic : children) {
            graphic.draw();
        }
    }
}
\end{lstlisting}

\subsection{Strategy}

\subsubsection{Назначение}

\textbf{Strategy} определяет семейство алгоритмов, инкапсулирует каждый из них и делает их взаимозаменяемыми.

\subsubsection{Минимальный пример}

\begin{lstlisting}[language=Java,caption={Strategy}]
interface SortStrategy {
    void sort(int[] array);
}

class BubbleSort implements SortStrategy {
    public void sort(int[] array) {
        System.out.println("Bubble sort");
    }
}

class QuickSort implements SortStrategy {
    public void sort(int[] array) {
        System.out.println("Quick sort");
    }
}

class Sorter {
    private SortStrategy strategy;

    public Sorter(SortStrategy strategy) {
        this.strategy = strategy;
    }

    public void sort(int[] array) {
        strategy.sort(array);
    }
}
\end{lstlisting}

\subsubsection{Алгоритм применения}

\begin{enumerate}
    \item Выделить изменяющийся алгоритм.
    \item Описать общий интерфейс стратегии.
    \item Создать конкретные стратегии.
    \item Передавать стратегию объекту-контексту.
    \item Вызывать алгоритм через интерфейс стратегии.
\end{enumerate}

\subsubsection{Иллюстрация}

\[
\texttt{Sorter} \rightarrow \texttt{SortStrategy}
\]

\[
\texttt{BubbleSort}, \quad \texttt{QuickSort}, \quad \texttt{MergeSort}.
\]

\subsection{Observer}

\subsubsection{Назначение}

\textbf{Observer} определяет зависимость ``один ко многим'', при которой изменение состояния одного объекта автоматически уведомляет зависимые объекты.

\subsubsection{Пример}

\[
\texttt{Subject} \rightarrow \texttt{Observer}_1, \texttt{Observer}_2, \dots, \texttt{Observer}_n.
\]

\subsubsection{Минимальный пример}

\begin{lstlisting}[language=Java,caption={Observer}]
import java.util.*;

interface Observer {
    void update(String message);
}

class Subject {
    private List<Observer> observers = new ArrayList<>();

    public void subscribe(Observer observer) {
        observers.add(observer);
    }

    public void notifyObservers(String message) {
        for (Observer observer : observers) {
            observer.update(message);
        }
    }
}

class EmailSubscriber implements Observer {
    public void update(String message) {
        System.out.println("Email: " + message);
    }
}

class SmsSubscriber implements Observer {
    public void update(String message) {
        System.out.println("SMS: " + message);
    }
}
\end{lstlisting}

\subsubsection{Алгоритм уведомления}

\begin{enumerate}
    \item Субъект хранит список наблюдателей.
    \item Наблюдатель подписывается на субъект.
    \item Состояние субъекта изменяется.
    \item Субъект вызывает метод уведомления у каждого наблюдателя.
    \item Каждый наблюдатель реагирует на изменение.
\end{enumerate}

\subsection{Command}

\subsubsection{Назначение}

\textbf{Command} превращает запрос в объект.

Это позволяет:
\begin{itemize}
    \item параметризовать объекты действиями;
    \item ставить команды в очередь;
    \item логировать команды;
    \item реализовывать отмену операций.
\end{itemize}

\subsubsection{Минимальный пример}

\begin{lstlisting}[language=Java,caption={Command}]
interface Command {
    void execute();
}

class Light {
    public void on() {
        System.out.println("Light on");
    }

    public void off() {
        System.out.println("Light off");
    }
}

class LightOnCommand implements Command {
    private Light light;

    public LightOnCommand(Light light) {
        this.light = light;
    }

    public void execute() {
        light.on();
    }
}

class RemoteControl {
    private Command command;

    public void setCommand(Command command) {
        this.command = command;
    }

    public void pressButton() {
        command.execute();
    }
}
\end{lstlisting}

\subsection{State}

\subsubsection{Назначение}

\textbf{State} позволяет объекту изменять своё поведение при изменении внутреннего состояния.

\subsubsection{Минимальный пример}

\begin{lstlisting}[language=Java,caption={State}]
interface State {
    void handle();
}

class DraftState implements State {
    public void handle() {
        System.out.println("Document is draft");
    }
}

class PublishedState implements State {
    public void handle() {
        System.out.println("Document is published");
    }
}

class Document {
    private State state;

    public void setState(State state) {
        this.state = state;
    }

    public void request() {
        state.handle();
    }
}
\end{lstlisting}

\subsection{Template Method}

\subsubsection{Назначение}

\textbf{Template Method} задаёт скелет алгоритма в базовом классе, позволяя подклассам переопределять отдельные шаги алгоритма.

\subsubsection{Минимальный пример}

\begin{lstlisting}[language=Java,caption={Template Method}]
abstract class DataProcessor {
    public final void process() {
        readData();
        transformData();
        saveData();
    }

    protected abstract void readData();
    protected abstract void transformData();

    protected void saveData() {
        System.out.println("Saving data");
    }
}

class CsvProcessor extends DataProcessor {
    protected void readData() {
        System.out.println("Reading CSV");
    }

    protected void transformData() {
        System.out.println("Transforming CSV");
    }
}
\end{lstlisting}

\subsubsection{Шаги алгоритма}

\begin{enumerate}
    \item Базовый класс определяет общий порядок действий.
    \item Некоторые шаги реализуются в базовом классе.
    \item Некоторые шаги объявляются абстрактными.
    \item Подклассы задают конкретную реализацию отдельных шагов.
\end{enumerate}

\subsection{Iterator}

\subsubsection{Назначение}

\textbf{Iterator} предоставляет способ последовательного доступа к элементам агрегированного объекта без раскрытия его внутреннего представления.

Этот шаблон лежит в основе стандартных итераторов коллекций.

\subsubsection{Минимальный пример}

\begin{lstlisting}[language=Java,caption={Iterator}]
interface MyIterator<T> {
    boolean hasNext();
    T next();
}

class ArrayIterator<T> implements MyIterator<T> {
    private T[] array;
    private int index = 0;

    public ArrayIterator(T[] array) {
        this.array = array;
    }

    public boolean hasNext() {
        return index < array.length;
    }

    public T next() {
        return array[index++];
    }
}
\end{lstlisting}

\section{Принципы проектирования ОО-систем}

\subsection{SOLID}

SOLID --- набор принципов объектно-ориентированного проектирования:
\begin{enumerate}
    \item Single Responsibility Principle;
    \item Open/Closed Principle;
    \item Liskov Substitution Principle;
    \item Interface Segregation Principle;
    \item Dependency Inversion Principle.
\end{enumerate}

\subsection{Принцип единственной ответственности}

\textbf{Single Responsibility Principle}:
\[
\text{у класса должна быть только одна причина для изменения.}
\]

Если класс одновременно отвечает за бизнес-логику, сохранение в базу данных и вывод на экран, он нарушает этот принцип.

Плохой пример:
\begin{lstlisting}[language=Java]
class Report {
    public void calculate() {}
    public void saveToFile() {}
    public void print() {}
}
\end{lstlisting}

Лучше разделить ответственности:
\begin{lstlisting}[language=Java]
class Report {
    public void calculate() {}
}

class ReportSaver {
    public void saveToFile(Report report) {}
}

class ReportPrinter {
    public void print(Report report) {}
}
\end{lstlisting}

\subsection{Принцип открытости/закрытости}

\textbf{Open/Closed Principle}:
\[
\text{программные сущности должны быть открыты для расширения, но закрыты для изменения.}
\]

Идея: новую функциональность лучше добавлять через новые классы, а не через постоянное изменение старого кода.

Например, вместо большого оператора \texttt{switch} по типам фигур лучше использовать полиморфизм:

\begin{lstlisting}[language=Java]
interface Shape {
    double area();
}
\end{lstlisting}

Новая фигура добавляется новым классом, не меняя код вычисления суммы площадей.

\subsection{Принцип подстановки Лисков}

\textbf{Liskov Substitution Principle}:
\[
\text{объекты подкласса должны быть заменяемы объектами базового класса без нарушения корректности программы.}
\]

Если функция ожидает объект типа \texttt{Shape}, то передача объекта \texttt{Circle} не должна ломать её логику.

Нарушение принципа часто возникает, когда подкласс усиливает предусловия или ослабляет постусловия методов базового класса.

\subsection{Принцип разделения интерфейсов}

\textbf{Interface Segregation Principle}:
\[
\text{клиенты не должны зависеть от методов, которые они не используют.}
\]

Плохой пример:
\begin{lstlisting}[language=Java]
interface Worker {
    void work();
    void eat();
}

class Robot implements Worker {
    public void work() {}
    public void eat() {
        throw new UnsupportedOperationException();
    }
}
\end{lstlisting}

Лучше:
\begin{lstlisting}[language=Java]
interface Workable {
    void work();
}

interface Eatable {
    void eat();
}
\end{lstlisting}

\subsection{Принцип инверсии зависимостей}

\textbf{Dependency Inversion Principle}:
\[
\text{модули верхнего уровня не должны зависеть от модулей нижнего уровня;}
\]
\[
\text{оба должны зависеть от абстракций.}
\]

Плохой пример:
\begin{lstlisting}[language=Java]
class MySQLDatabase {
    void save(String data) {}
}

class UserService {
    private MySQLDatabase database = new MySQLDatabase();
}
\end{lstlisting}

Лучше:
\begin{lstlisting}[language=Java]
interface Database {
    void save(String data);
}

class MySQLDatabase implements Database {
    public void save(String data) {}
}

class UserService {
    private Database database;

    public UserService(Database database) {
        this.database = database;
    }
}
\end{lstlisting}

\section{Связность и сцепление}

\subsection{Связность}

\textbf{Связность} модуля показывает, насколько тесно элементы внутри одного модуля связаны между собой.

Высокая связность обычно является хорошим свойством.

Класс должен иметь чёткую ответственность, а его методы должны работать на общую цель.

\subsection{Сцепление}

\textbf{Сцепление} показывает степень зависимости одного модуля от другого.

Низкое сцепление обычно является хорошим свойством.

ООП и шаблоны проектирования часто направлены на:
\[
\text{высокую связность} + \text{низкое сцепление}.
\]

\section{Сравнение наследования и композиции}

\begin{longtable}{|p{5cm}|p{5cm}|p{5cm}|}
\hline
\textbf{Критерий} & \textbf{Наследование} & \textbf{Композиция} \\
\hline
Отношение & is-a & has-a \\
\hline
Связь классов & Более жёсткая & Более гибкая \\
\hline
Повторное использование & Через базовый класс & Через включение объекта \\
\hline
Изменение поведения & Обычно через переопределение & Через замену компонента \\
\hline
Риск нарушения инкапсуляции & Выше & Ниже \\
\hline
Пример & \texttt{Circle extends Shape} & \texttt{Car has Engine} \\
\hline
\end{longtable}

\section{Исключения как объектный механизм}

\subsection{Понятие исключения}

\textbf{Исключение} --- объект, описывающий ошибочную или нестандартную ситуацию во время выполнения программы.

Исключения позволяют отделить обычную логику программы от обработки ошибок.

\begin{lstlisting}[language=Java,caption={Исключения}]
public class Main {
    public static int divide(int a, int b) {
        if (b == 0) {
            throw new IllegalArgumentException("Division by zero");
        }
        return a / b;
    }

    public static void main(String[] args) {
        try {
            System.out.println(divide(10, 0));
        } catch (IllegalArgumentException e) {
            System.out.println(e.getMessage());
        }
    }
}
\end{lstlisting}

\section{Сериализация объектов}

\subsection{Понятие сериализации}

\textbf{Сериализация} --- преобразование объекта в последовательность байтов или текстовый формат для хранения или передачи.

\textbf{Десериализация} --- восстановление объекта из сохранённого представления.

Применения:
\begin{itemize}
    \item сохранение состояния программы;
    \item передача объектов по сети;
    \item кеширование;
    \item обмен данными между сервисами.
\end{itemize}

\section{Краткая сравнительная таблица шаблонов}

\begin{longtable}{|p{4cm}|p{4cm}|p{6cm}|}
\hline
\textbf{Шаблон} & \textbf{Группа} & \textbf{Основная идея} \\
\hline
Singleton & Порождающий & Единственный экземпляр класса \\
\hline
Factory Method & Порождающий & Делегирование создания объектов подклассам \\
\hline
Abstract Factory & Порождающий & Создание семейств связанных объектов \\
\hline
Builder & Порождающий & Пошаговое создание сложного объекта \\
\hline
Prototype & Порождающий & Создание объекта копированием прототипа \\
\hline
Adapter & Структурный & Совмещение несовместимых интерфейсов \\
\hline
Decorator & Структурный & Динамическое добавление поведения \\
\hline
Facade & Структурный & Простой интерфейс к сложной подсистеме \\
\hline
Proxy & Структурный & Заместитель объекта с контролем доступа \\
\hline
Composite & Структурный & Единообразная работа с объектами и их группами \\
\hline
Strategy & Поведенческий & Взаимозаменяемые алгоритмы \\
\hline
Observer & Поведенческий & Уведомление зависимых объектов об изменениях \\
\hline
Command & Поведенческий & Представление запроса в виде объекта \\
\hline
State & Поведенческий & Изменение поведения при изменении состояния \\
\hline
Template Method & Поведенческий & Скелет алгоритма в базовом классе \\
\hline
Iterator & Поведенческий & Последовательный обход контейнера \\
\hline
\end{longtable}

\section{Минимальный пример совместного применения принципов ООП}

Рассмотрим задачу: нужно отправлять уведомления разными способами.

Плохой вариант:

\begin{lstlisting}[language=Java]
class NotificationService {
    public void send(String type, String message) {
        if (type.equals("email")) {
            System.out.println("Email: " + message);
        } else if (type.equals("sms")) {
            System.out.println("SMS: " + message);
        }
    }
}
\end{lstlisting}

Недостатки:
\begin{itemize}
    \item при добавлении нового способа уведомления нужно изменять класс;
    \item нарушается принцип открытости/закрытости;
    \item логика разных уведомлений смешана.
\end{itemize}

Хороший вариант:

\begin{lstlisting}[language=Java]
interface Notifier {
    void send(String message);
}

class EmailNotifier implements Notifier {
    public void send(String message) {
        System.out.println("Email: " + message);
    }
}

class SmsNotifier implements Notifier {
    public void send(String message) {
        System.out.println("SMS: " + message);
    }
}

class NotificationService {
    private Notifier notifier;

    public NotificationService(Notifier notifier) {
        this.notifier = notifier;
    }

    public void notify(String message) {
        notifier.send(message);
    }
}
\end{lstlisting}

Здесь используются:
\begin{itemize}
    \item абстракция --- интерфейс \texttt{Notifier};
    \item инкапсуляция --- конкретные способы отправки скрыты в классах;
    \item полиморфизм --- сервис работает с любым объектом \texttt{Notifier};
    \item композиция --- сервис содержит объект уведомления;
    \item Strategy --- способ отправки сообщения является заменяемой стратегией;
    \item Dependency Inversion --- сервис зависит от интерфейса, а не от конкретного класса.
\end{itemize}

\section{Выводы}

Объектно-ориентированное программирование строится вокруг понятий класса, объекта, состояния, поведения и взаимодействия объектов.

Основные принципы ООП:
\begin{itemize}
    \item абстракция позволяет выделять главное;
    \item инкапсуляция защищает внутреннее состояние объекта;
    \item наследование позволяет строить иерархии и повторно использовать код;
    \item полиморфизм обеспечивает единообразную работу с разными типами объектов.
\end{itemize}

Интерфейсы, библиотеки классов, контейнеры, итераторы, рефлексия и корутины образуют важную часть объектного окружения современных языков программирования.

Шаблоны проектирования представляют собой проверенные решения типичных задач проектирования. Они делятся на порождающие, структурные и поведенческие. Правильное применение шаблонов повышает гибкость, расширяемость и сопровождаемость программ, однако чрезмерное использование шаблонов может сделать архитектуру излишне сложной.

\end{document}
